%% Sometimes when a section can't be nicely modelled with the \entry[]... mechanism; hack our own
\makerubrichead{Publications and Patents}

\vspace{0.1cm}
\begin{center}
\begin{tabular}{|l|l|}
\hline
\textbf{Metric} & \textbf{Value} \\
\hline
N. journal papers & $15$ \\
N. conference papers & $23$ \\
N. patents & $2$ ($1$ pending) \\
N. citations (Google Scholar) & $530+$ \\
h-index (Google Scholar) & $13$ \\
\hline
\end{tabular}
\end{center}

% #1: label prefix, #2: keyword assigned to the source .bib file (see main_cv.tex), #3: title
\newcommand{\printpublist}[3]{%
	\newrefcontext[labelprefix=#1]%
	\printbibliography[keyword=#2,heading={subbibliography},title={#3},resetnumbers=true]%
}

\vspace{0.1cm}

\printpublist{J}{cvjournal}{Peer-Reviewed Journal Papers}
\printpublist{JR}{cvjournalrev}{Peer-Reviewed Journal Papers \textit{Under Review}}

\printpublist{C}{cvconf}{Peer-Reviewed Conference Papers}
% \printpublist{CR}{cvconfrev}{Peer-Reviewed Conference Papers \textit{Under Review}}

\printpublist{W}{cvworkshop}{Peer-Reviewed Workshop Papers}
\printpublist{O}{cvother}{Other Publications}
% \printbibliographyfromfile{invited_talks.bib}{Invited Talks}
\printpublist{PS}{cvposter}{Poster Sessions}

\printpublist{PT}{cvpatent}{Patents}
\endrefcontext
